% ============================================================
% Script de soutenance — Sheryne (diapos 3, 6 à 9) — version normale
% Compilation : pdflatex script_sheryne.tex
% ============================================================
\documentclass[11pt,a4paper]{article}

\usepackage[utf8]{inputenc}
\usepackage[T1]{fontenc}
\usepackage[french]{babel}
\usepackage[margin=2.4cm]{geometry}
\usepackage{amsmath,amssymb}
\usepackage{booktabs}
\usepackage{parskip}

\newcommand{\diapo}[2]{\section*{Diapo #1 --- #2}}
\newenvironment{comprendre}%
{\par\medskip\noindent\hrulefill\par\nopagebreak\textbf{Pour bien comprendre :}\itshape\par}%
{\par\nopagebreak\noindent\hrulefill\par\medskip}

\title{\textbf{Script de soutenance --- Sheryne}\\
	\large Effet tunnel et temps de traversée d'une barrière de potentiel\\
	\normalsize Diapos 3, 6--9 --- environ 5 minutes}
\author{}
\date{Soutenance juin 2026 --- oraux du 22 au 26 juin}

\begin{document}
	\maketitle
	
	\subsection*{Répartition des rôles}
	
	\begin{center}
		\begin{tabular}{llc}
			\toprule
			Qui & Diapos & Durée visée \\
			\midrule
			Myriam & 1, 2, 4--5 : problématique, paquet d'ondes, états stationnaires & $\approx$ 5 min \\
			\textbf{Sheryne} & 3, 6--9 : motivation, méthode numérique, validation, premiers résultats & $\approx$ 5 min \\
			Mattéo & 10--13 : effet Hartman, études paramétriques, conclusion & $\approx$ 5,5 min \\
			\bottomrule
		\end{tabular}
	\end{center}
	
	Total : $\approx$ 15,5 minutes, puis 10 minutes de questions. Le texte ci-dessous est à dire
	à l'oral ; les phrases \textbf{en gras} sont les messages essentiels, à ne jamais
	sauter même si le temps presse.
	
	% ------------------------------------------------------------
	\diapo{3}{Motivation : de l'onde plane au paquet d'ondes (30 secondes)}
	
	Merci Myriam. Avant le paquet gaussien, un mot sur d'où il vient.
	
	Une onde plane seule, $\Psi_0\,e^{i(kx-\omega t)}$, n'est pas physiquement
	acceptable : son module au carré est constant sur tout l'espace, donc
	\textbf{non normalisable}. Pour décrire une particule réellement localisée, on
	superpose plusieurs ondes planes de vecteurs d'onde voisins. Les interférences
	entre elles créent une enveloppe qui localise la particule : c'est le
	battement que vous voyez sur la figure. En passant à une superposition
	continue sur tous les $k$, on obtient le paquet d'ondes gaussien.
	
	Je rends la parole à Myriam pour le détailler.
	
	% ------------------------------------------------------------
	\diapo{6}{Méthode numérique : schéma de Crank--Nicolson (1,5 minute)}
	
	Merci Myriam, je reprends la suite. Les outils théoriques que vous venez de
	voir décrivent soit un paquet libre, soit une onde plane face à la barrière. Mais l'évolution complète
	d'un \emph{paquet} face à une barrière n'a pas de solution analytique simple :
	\textbf{il faut résoudre numériquement l'équation de Schrödinger dépendante du
		temps}.
	
	Nous avons choisi le schéma de Crank--Nicolson. Son principe : au lieu d'évaluer
	le membre de droite de l'équation soit à l'instant présent, soit à l'instant
	suivant, \textbf{on prend la moyenne des deux instants}. Cela donne l'équation
	discrétisée à l'écran, où le coefficient $r$ contient le pas de temps et le pas
	d'espace, et où $s_i$ contient le potentiel au point $x_i$. Comme notre potentiel
	varie dans l'espace --- il est nul partout sauf sur la barrière --- la diagonale
	du système varie d'un point à l'autre. \textbf{Et c'est exactement ce que fait
		le code à droite} : la ligne \texttt{s\_int} porte cette variation, et c'est
	elle qui rend la diagonale du système différente à chaque point.
	
	Pourquoi ce schéma plutôt qu'un schéma explicite plus simple ? Trois raisons.
	\textbf{Un : il est unitaire} --- la norme de la fonction d'onde, c'est-à-dire la
	probabilité totale, est conservée à $10^{-6}$ près sur toute la simulation. Pour
	de la mécanique quantique, c'est la propriété essentielle. \textbf{Deux : il est
		inconditionnellement stable}, quel que soit le pas de temps, et précis à l'ordre
	deux en temps et en espace. \textbf{Trois : il reste rapide}, car le système à
	résoudre à chaque pas de temps est tridiagonal, et la méthode de Thomas le résout
	en un nombre d'opérations proportionnel au nombre de points --- vous la voyez
	appelée tout en bas du code, juste après avoir construit la diagonale.
	
	Nos paramètres de référence : une grille de 1500 points sur $[-30, 30]$, 3000 pas
	de temps jusqu'à $t = 6$ ; un paquet de vitesse $k_0 = 5$, donc d'énergie
	$E = 12{,}5$, parti de $x_0 = -10$ ; et une barrière de largeur $a = 1$ et de
	hauteur $V_0 = 15$, donc \textbf{plus haute que l'énergie du paquet} : on est
	bien en régime tunnel.
	
	\begin{comprendre}
		« Unitaire » veut dire que le schéma multiplie la fonction d'onde par un
		opérateur de norme 1, comme le fait l'évolution quantique exacte. Un schéma
		d'Euler explicite, lui, amplifie la norme à chaque pas : la probabilité totale
		dépasse 1 et la simulation explose. C'est la réponse à donner si le jury demande
		« pourquoi pas Euler ? ». « Tridiagonal » signifie que chaque ligne du système ne
		couple un point qu'à ses deux voisins immédiats --- conséquence directe de la
		discrétisation de la dérivée seconde.
	\end{comprendre}
	
	% ------------------------------------------------------------
	\diapo{7}{Validation : particule libre (1 minute)}
	
	Avant d'allumer la barrière, on valide le solveur sur un cas dont on connaît la
	réponse : la particule libre, $V_0 = 0$.
	
	Notre définition opérationnelle du temps de traversée : \textbf{on suit le pic de
		$|\psi|^2$, et on note les instants où il passe en $x = 0$ puis en $x = a$}. La
	différence donne le temps de traversée numérique $\tau_0$, à comparer à la
	prédiction évidente : la distance divisée par la vitesse de groupe, soit
	$a / v_g = 0{,}2$. C'est exactement ce que fait la fonction à l'écran,
	\texttt{temps\_pic\_en} : elle repère où l'amplitude est maximale et renvoie
	l'instant correspondant --- on l'appelle une fois à l'entrée, une fois à la
	sortie.
	
	Le tableau montre l'accord : entrée mesurée à $t = 1{,}9987$ pour $2$ attendu,
	sortie à $2{,}1907$ pour $2{,}2$, soit \textbf{$\tau_0 = 0{,}192$ contre
		$0{,}200$ théorique : un écart de 4\,\%}, qui s'explique par la résolution
	temporelle de la grille et par l'étalement du paquet pendant le vol. Et la norme
	reste égale à 1 à $10^{-6}$ près sur les 3000 pas.
	
	\textbf{Le solveur est validé : on peut maintenant allumer la barrière.}
	
	\begin{comprendre}
		Cette étape de validation est méthodologiquement importante : toute mesure faite
		ensuite avec la barrière héritera de cette incertitude de l'ordre de 4\,\%. Quand
		Mattéo comparera $\tau_t$ numérique et théorique, il faudra garder cet ordre de
		grandeur en tête avant de crier à la découverte.
	\end{comprendre}
	
	% ------------------------------------------------------------
	\diapo{8}{Évolution face à la barrière (1 minute)}
	
	Voici maintenant le film de la simulation complète, en quatre instants.
	
	Le paquet arrive de la gauche, heurte la barrière --- la zone grisée --- et se
	sépare en deux : \textbf{un paquet réfléchi, largement dominant, qui repart vers
		la gauche, et un petit paquet transmis qui continue vers la droite}. Ce petit
	paquet transmis, c'est l'effet tunnel en images : \textbf{classiquement, avec
		$E < V_0$, la particule serait réfléchie avec certitude} et il n'y aurait
	strictement rien à droite.
	
	En bleu, nous avons superposé le paquet libre, celui de la diapo précédente : il
	avance à la même vitesse de groupe, ce qui permet de comparer les positions ---
	et c'est cette comparaison qui servira à mesurer le retard, ou l'avance, du
	paquet transmis.
	
	\begin{comprendre}
		Pendant l'impact, on voit des franges devant la barrière : c'est l'interférence
		entre le paquet incident et le paquet réfléchi, qui coexistent au même endroit au
		même moment. C'est un bon détail à mentionner si le jury demande de commenter la
		figure.
	\end{comprendre}
	
	% ------------------------------------------------------------
	\diapo{9}{Probabilités de réflexion et de transmission (1 minute)}
	
	Pour quantifier, on intègre $|\psi|^2$ de chaque côté de la barrière au cours du
	temps : c'est le code à droite, avec un masque qui sélectionne la zone après
	la barrière, et une ligne qui intègre la densité dessus. La courbe du haut le
	confirme : \textbf{la norme totale reste constante}, c'est l'unitarité de
	Crank--Nicolson à l'\oe uvre.
	
	À la fin de la simulation, la probabilité d'être passé à droite --- c'est
	exactement ce que calcule cette ligne --- vaut $8{,}1\times10^{-2}$. Or la formule
	des états stationnaires de Myriam prédit
	$|T(k_0)|^2 = 2{,}5\times10^{-2}$ pour une onde plane de vecteur d'onde $k_0$.
	\textbf{Trois fois plus que la théorie : est-ce une erreur ? Non --- et c'est un
		de nos résultats les plus intéressants.}
	
	L'explication : notre paquet n'est pas une onde plane. Il contient toute une
	gamme de vecteurs d'onde autour de $k_0$, et la transmission $|T|^2$ croît très
	vite avec $k$. \textbf{La barrière favorise donc fortement les composantes
		rapides du paquet : elle agit comme un filtre passe-haut en vitesse.} Ce
	filtrage spectral est la clé de toute la suite : c'est lui qui explique, comme
	vous le verrez avec Mattéo, l'écart entre le temps de traversée mesuré et la
	théorie de Hartman.
	
	Mattéo, à toi.
	
	\begin{comprendre}
		Pour fixer l'image : la barrière reçoit un mélange de composantes lentes et
		rapides, et ne laisse passer presque que les rapides. Le paquet transmis est donc
		plus rapide \emph{en moyenne} que le paquet incident. C'est pour cela que la
		probabilité mesurée dépasse la prédiction « onde plane à $k_0$ » : la moyenne de
		$|T(k)|^2$ sur le spectre du paquet est dominée par les grands $k$, bien mieux
		transmis.
	\end{comprendre}
	
	% ------------------------------------------------------------	
\end{document}